% TOP_PROBLEM: {"id": 6600015, "title": "L. Sadun — Problem", "category_id": 6, "category": {"id": 6, "name": "geometry", "display_name": "Geometry", "description": "Euclidean and non-Euclidean geometry, geometric structures.", "slug": "geometry", "order_index": 6, "created_at": "2026-07-31T15:26:25.671Z"}, "status": "open"}
% FIRST_POSTED: null
% ENTRY_KIND: statement_only
% Imported from: batch13.tex; UnsolvedMath ID 6600015
\documentclass[11pt]{article}
\usepackage[margin=1in]{geometry}
\usepackage{amsmath,amssymb,mathtools}
\usepackage{fontspec,unicode-math}
\setmainfont{Latin Modern Roman}
\setsansfont{Latin Modern Sans}
\setmathfont{Latin Modern Math}
\usepackage{xurl,hyperref}
\hypersetup{colorlinks=true,urlcolor=blue,linkcolor=blue}
\newcommand{\sref}[2]{\hyperlink{source:#1:#2}{[S#2]}}
\newcommand{\cataloguescope}[1]{\paragraph{Recorded mathematical question.}#1}
\begin{document}
% UnsolvedMath ID 6600015; source catalogue code AMR-065-0015
\setcounter{section}{6600014}
\section{L. Sadun — Problem}\label{6600015}
{\small\sffamily UnsolvedMath ID 6600015 · Geometry}
\subsection{Definitions and mathematical statement}

\paragraph{Shared notation.}$\mathbb N=\{1,2,\ldots\}$, and $\mathbb Z,\mathbb Q,\mathbb R,\mathbb C$ denote the integers, rationals, reals and complex numbers. Any other conventions are specified locally.


For point patterns associated to tilings by a fixed reference point in each tile, BD equivalence to a lattice means a bijection with uniformly bounded displacement; BL equivalence means a bijection with uniformly bounded two-sided distance distortion. Linear repetitivity (LR) means that some $C$ makes every radius-$r$ patch occur in every ball of radius $Cr$, for all $r\geq1$. A Meyer point pattern is Delone and has uniformly discrete difference set. Pure point diffraction means that its diffraction measure is a sum of point masses. The tiling hull is the closure of translates in the local tiling topology. Mutual local derivability (MLD) means that each tiling can be recovered from the other by translation-compatible rules using a uniformly bounded neighborhood. Topological conjugacy means a homeomorphism of hulls commuting with translations; homeomorphism alone means a homeomorphism of hulls without that equivariance requirement.
\paragraph{Statement recorded in the supplied source.}
For geometric properties such as BD or BL equivalence to a lattice, LR, pure point diffraction, or the Meyer property, classify the tilings for which that property is also possessed by every tiling equivalent to them under MLD, under topological conjugacy, or under homeomorphism of tiling hulls. Treat each property and each equivalence relation as its own classification question. \sref{6600015}{2}
\subsection{Short English statement}
Classify tilings whose chosen geometric property survives throughout their entire MLD, conjugacy, or hull-homeomorphism class.
\subsection{Sources}
\begin{description}
\item[\hypertarget{source:6600015:1}{S1}] ULAM.AI and the UnsolvedMath Contributors, \emph{UnsolvedMath: A Curated Collection of Open Mathematics Problems}, record 6600015 (AMR-065-0015). CC BY 4.0. \url{https://huggingface.co/datasets/ulamai/UnsolvedMath}.
\item[\hypertarget{source:6600015:2}{S2}] Faustin Adiceam (compiler), \emph{Open Problems and Conjectures related to the Theory of Mathematical Quasicrystals}, arXiv:1604.06280v2 (2016), Introduction, pp. 1--2; Section 2.7.1, Problem 2.7.1, p. 9. \url{https://arxiv.org/pdf/1604.06280}.
\end{description}
\label{end:6600015}
\end{document}
